\documentclass[11pt]{article}
\usepackage[a4paper,margin=1in]{geometry}
\usepackage{amsmath,amssymb,mathtools,bm}
\usepackage{enumitem,booktabs,array}
\usepackage{tcolorbox}
\usepackage{hyperref}
\usepackage[protrusion=true,expansion=false]{microtype}
\hypersetup{colorlinks=true,linkcolor=blue,urlcolor=blue,citecolor=blue}
\setlist[itemize]{topsep=2pt,itemsep=2pt}
\setlist[enumerate]{topsep=2pt,itemsep=2pt}
\newtcolorbox{keybox}{colback=blue!3!white,colframe=blue!40!black,boxrule=0.4pt,arc=1mm}
\newtcolorbox{alertbox}{colback=red!3!white,colframe=red!40!black,boxrule=0.4pt,arc=1mm}
\newtcolorbox{answerbox}{colback=green!3!white,colframe=green!40!black,boxrule=0.4pt,arc=1mm}
\newtcolorbox{drillbox}{colback=yellow!5!white,colframe=orange!60!black,boxrule=0.4pt,arc=1mm}
\newcommand{\EV}{\mathbb{E}}
\newcommand{\Var}{\mathrm{Var}}
\newcommand{\Cov}{\mathrm{Cov}}
\newcommand{\blank}[1]{\underline{\hspace{#1}}}

\title{W13-02: Addoum, Ng \& Ortiz-Bobea (2020, RFS) \\
\large ``Temperature Shocks and Establishment Sales'' --- Active Recall Workbook}
\author{Banking \& Risk Management --- Exam Prep}
\date{\today}
\begin{document}\maketitle

\begin{keybox}
\textbf{Paper in one sentence.} Using county-level panel data on $\sim$900{,}000 US establishments 1990--2015 merged with daily temperature records, ANO show that hot days ($>90^\circ$F) reduce establishment sales in exposed industries by $\sim$0.2--0.5\% per extra hot day --- firm-level productivity evidence that physical climate damages are already materialising.

\textbf{Placement.} Canonical climate-and-real-activity paper. Empirically quantifies physical-risk damage at the firm level, complementing the climate-pricing literature (BGL, BBGL, MS) with real-outcome estimates.
\end{keybox}

\section*{Round 1: Complete Walk-Through}

\subsection*{1.1 Data}
Matched establishment-level sales from NETS (National Establishment Time Series) with county-level daily temperature records from PRISM. Industries classified by heat vulnerability (construction, agriculture, tourism more exposed; finance, professional services less).

\subsection*{1.2 Empirical spec}
Panel regression with establishment and year FE:
\[
\Delta\ln\text{Sales}_{it}=\alpha_i+\delta_t+\sum_k\beta_k\cdot\text{DaysInBin}_{k,ct}+X_{it}'\gamma+\varepsilon_{it},
\]
where DaysInBin$_{k,ct}$ is the number of days in county $c$, year $t$ in temperature bin $k$ (e.g., $80$--$85$F, $85$--$90$F, $>90$F). Reference bin is a moderate temperature range.

\subsection*{1.3 Headline results}
\begin{itemize}
\item Each extra day above 90$^\circ$F reduces establishment sales by $\sim$0.2--0.5\% in exposed industries.
\item Effect is concave (very hot days much worse than moderately hot days).
\item Heterogeneity: agriculture, construction, tourism, retail most exposed; finance and professional services not.
\item Aggregate: US establishment sales roughly 1\% lower per year of high-heat exposure.
\end{itemize}

\subsection*{1.4 Interpretation}
\begin{itemize}
\item Physical climate damages are firm-level and measurable.
\item Adaptation is partial: A/C investment, shift scheduling, reducing outdoor work.
\item Implications for forward-looking firm valuations and cost of capital.
\end{itemize}

\subsection*{1.5 Caveats}
\begin{alertbox}
\begin{itemize}
\item NETS data quality: sales are sometimes imputed.
\item Climate change prediction would require extrapolating from observed weather shocks.
\item Adaptation could reduce damages over longer horizons.
\end{itemize}
\end{alertbox}

\section*{Round 2: Level 1 Fill-in}
\begin{enumerate}
\item Sample: $\sim$\blank{2cm} US establishments, \blank{1cm}--\blank{1cm}.
\item Weather data source: \blank{1.5cm}.
\item Per-hot-day effect: $\sim$\blank{1cm}--\blank{1cm}\%.
\item Most exposed industries: \blank{2cm}, construction, tourism.
\item Least exposed: \blank{1cm}.
\end{enumerate}

\section*{Round 3: Level 2 Prompts}
\begin{enumerate}
\item Write the binned-temperature specification and explain identification.
\item Why is the effect concave in temperature?
\item Discuss adaptation margins that firms might take.
\item How do these results translate into long-run valuations?
\end{enumerate}

\section*{Round 4: Minimal Scaffolding}
From a blank page: (i) data, (ii) spec, (iii) magnitudes, (iv) heterogeneity, (v) caveats.

\section*{Round 5: Blank Page}
Essay: physical climate risk and firm-level real activity.

\vspace{6pt}
\begin{drillbox}
\textbf{Speed Drill.} (1) Sample size. (2) Per-hot-day magnitude. (3) Exposed industries. (4) Weather data. (5) Concavity finding.
\end{drillbox}

\section*{Quick Reference \& Answer Key}
\begin{answerbox}
\begin{itemize}
\item \textbf{Data.} 900k establishments, 1990--2015, with PRISM weather.
\item \textbf{Effect.} $\sim$0.2--0.5\% sales loss per extra day $>90^\circ$F.
\item \textbf{Concavity.} Very hot days worst.
\item \textbf{Industries.} Agriculture, construction, tourism most exposed.
\end{itemize}
\end{answerbox}

\begin{keybox}
\textbf{30-second exam pitch.} Addoum, Ng \& Ortiz-Bobea (2020) document physical climate-risk damages at the firm level. Merging $\sim$900{,}000 US establishment records from NETS with county-level daily temperature from PRISM, 1990--2015, they estimate binned-temperature sales regressions with establishment and year FE. Each additional day above 90$^\circ$F cuts establishment sales by $\sim$0.2--0.5\% in exposed industries (agriculture, construction, tourism, retail); finance and professional services are largely unaffected. Effects are concave: extremely hot days are disproportionately damaging. Aggregating, US establishment sales are roughly 1\% lower per year of high-heat exposure. The paper is the canonical micro-level quantification of physical climate risk, showing that damages are already visible and that adaptation is only partial.
\end{keybox}

\end{document}
